%%
%% part5-examples/ch28-programming-book.tex
%%
%% Chapter 28: A Programming Book — A complete sample
%% chapter from a programming textbook, focusing on code,
%% algorithms, and data structures.
%%

\chapter{کتاب برنامه‌نویسی}
\label{chap:programming-book}

\section{چرا این موضوع مهم است؟}
\label{sec:programming-book-why}

در فصل قبل، یک فصل نمونه از کتاب ریاضی
دیدیم که بر قضیه، اثبات و معادلات تمرکز
داشت. در این فصل، یک فصل نمونه از کتاب
برنامه‌نویسی را می‌بینیم که بر کد، الگوریتم
و ساختار داده تمرکز دارد.

هدف این فصل، نشان دادن این است که چگونه
می‌توان از قابلیت‌های \lr{MatinBook} برای
نوشتن کتاب‌های برنامه‌نویسی استفاده کرد.

\mnote{کتاب‌های برنامه‌نویسی چالش‌های خاص خود را دارند: کد، الگوریتم، و توضیحات فارسی باید در کنار هم هماهنگ باشند.}

\section{ساختار فصل نمونه}
\label{sec:programming-book-structure}

فصل نمونه، از یک کتاب برنامه‌نویسی خیالی
با عنوان «مقدمه‌ای بر ساختمان داده‌ها»
است. این فصل شامل بخش‌های زیر است:

\begin{itemize}
    \item مقدمه.
    \item کد نمونه (پایتون و \lr{C++}).
    \item الگوریتم‌ها.
    \item جدول مقایسه.
    \item تمرین‌ها.
    \item راه‌حل‌ها.
\end{itemize}

\section{مقدمه}
\label{sec:programming-book-intro}

ساختمان داده\index{ساختمان داده}، روشی برای
سازمان‌دهی و ذخیره‌سازی داده‌ها در حافظه‌ی
کامپیوتر است. انتخاب ساختمان داده‌ی مناسب،
تأثیر مستقیمی بر کارایی الگوریتم‌ها دارد.

در این فصل، با ساختمان داده‌ی \textbf{پشته}
و کاربردهای آن آشنا می‌شویم.

\mnote{پشته یکی از ساده‌ترین و در عین حال پرکاربردترین ساختمان داده‌هاست.}

\section{پشته}
\label{sec:programming-book-stack}

\subsection{تعریف پشته}
\label{sec:programming-book-stack-def}

\begin{definition}[پشته]
    \textbf{پشته}\index{پشته} یک ساختمان
    داده‌ی خطی است که در آن درج و حذف
    عناصر از یک سر انجام می‌شود. این
    سر را \textbf{بالای پشته} می‌نامیم.
    \label{def:stack}
\end{definition}

طبق \cref{def:stack}، پشته از اصل
\textbf{\lr{LIFO}}\index{LIFO} (آخرین
ورودی، اولین خروجی) پیروی می‌کند.

\begin{example}
    عملیات اصلی پشته:
    \begin{itemize}
        \item \lr{push(x)}: درج \lr{x} در
        بالای پشته.
        \item \lr{pop()}: حذف و بازگشت عنصر
        بالای پشته.
        \item \lr{top()}: بازگشت عنصر بالای
        پشته بدون حذف آن.
        \item \lr{empty()}: بررسی خالی بودن
        پشته.
    \end{itemize}
    \label{ex:stack-operations}
\end{example}

\subsection{پیاده‌سازی در پایتون}
\label{sec:programming-book-python}

پیاده‌سازی پشته در پایتون با لیست:

\begin{latin}
\begin{minted}{python}
class Stack:
    """A simple stack implementation using a Python list."""

    def __init__(self) -> None:
        self._data: list = []

    def push(self, item) -> None:
        """Add an item to the top of the stack."""
        self._data.append(item)

    def pop(self):
        """Remove and return the top item."""
        if self.is_empty():
            raise IndexError("pop from empty stack")
        return self._data.pop()

    def peek(self):
        """Return the top item without removing it."""
        if self.is_empty():
            raise IndexError("peek from empty stack")
        return self._data[-1]

    def is_empty(self) -> bool:
        """Return True if the stack is empty."""
        return len(self._data) == 0

    def size(self) -> int:
        """Return the number of items in the stack."""
        return len(self._data)


# |\pc{استفاده از پشته}|
stack = Stack()
stack.push(10)
stack.push(20)
stack.push(30)
print(stack.pop())   # 30
print(stack.peek())  # 20
\end{minted}
\end{latin}
\codecaption{پیاده‌سازی پشته در پایتون}
\label{code:stack-python}

\mnote{در پایتون، لیست به‌عنوان ساختار پایه برای پشته استفاده می‌شود. عملیات \lr{append} و \lr{pop} هر دو \lr{$O(1)$} هستند.}

\subsection{پیاده‌سازی در \lr{C++}}
\label{sec:programming-book-cpp}

پیاده‌سازی پشته در \lr{C++} با قالب:

\begin{latin}
\begin{minted}{cpp}
#include <vector>
#include <stdexcept>

template <typename T>
class Stack {
private:
    std::vector<T> data;

public:
    void push(const T& item) {
        data.push_back(item);
    }

    T pop() {
        if (data.empty()) {
            throw std::runtime_error("pop from empty stack");
        }
        T item = data.back();
        data.pop_back();
        return item;
    }

    const T& peek() const {
        if (data.empty()) {
            throw std::runtime_error("peek from empty stack");
        }
        return data.back();
    }

    bool empty() const {
        return data.empty();
    }

    size_t size() const {
        return data.size();
    }
};
\end{minted}
\end{latin}
\codecaption{پیاده‌سازی پشته در \lr{C++}}
\label{code:stack-cpp}

\mnote{در \lr{C++}، از قالب (\lr{template}) استفاده می‌کنیم تا پشته با هر نوع داده‌ای کار کند.}

\section{الگوریتم‌ها}
\label{sec:programming-book-algorithms}

\subsection{بررسی توازن پرانتزها}
\label{sec:programming-book-parentheses}

یکی از کاربردهای کلاسیک پشته، بررسی
توازن پرانتزها در یک عبارت است.

\begin{latin}
\begin{algorithm}
\caption{بررسی توازن پرانتزها}
\label{alg:balanced-parentheses}
\begin{algorithmic}[1]
    \Require A string $s$ containing parentheses
    \Ensure \textbf{true} if the parentheses are balanced,
            \textbf{false} otherwise
    \State Create an empty stack $S$
    \For{each character $c$ in $s$}
        \If{$c$ is an opening bracket}
            \State $S.\text{push}(c)$
        \ElsIf{$c$ is a closing bracket}
            \If{$S$ is empty}
                \State \Return \textbf{false}
            \EndIf
            \State $top \gets S.\text{pop}()$
            \If{$top$ does not match $c$}
                \State \Return \textbf{false}
            \EndIf
        \EndIf
    \EndFor
    \State \Return $S$ is empty
\end{algorithmic}
\end{algorithm}
\end{latin}

\mnote{این الگوریتم در کامپایلرها، ویرایشگرهای متن و تحلیل‌گرهای نحوی کاربرد دارد.}

\subsection{پیچیدگی}
\label{sec:programming-book-complexity}

\begin{theorem}[پیچیدگی بررسی توازن]
    الگوریتم \cref{alg:balanced-parentheses}
    دارای پیچیدگی زمانی \lr{$O(n)$} و
    پیچیدگی فضایی \lr{$O(n)$} است، که
    \lr{$n$} طول رشته‌ی ورودی است.
    \label{thm:balanced-complexity}
\end{theorem}

\begin{proof}
    هر کاراکتر رشته دقیقاً یک بار پردازش
    می‌شود (حلقه‌ی \lr{for} در خط ۲). عملیات
    \lr{push} و \lr{pop} هر دو \lr{$O(1)$}
    هستند. بنابراین پیچیدگی زمانی
    \lr{$O(n)$} است.

    در بدترین حالت (رشته‌ای با \lr{$n$}
    پرانتز باز)، پشته دارای \lr{$n$} عنصر
    خواهد بود. بنابراین پیچیدگی فضایی
    \lr{$O(n)$} است.
\end{proof}

\section{جدول مقایسه}
\label{sec:programming-book-comparison}

جدول \cref{tab:stack-comparison} پیاده‌سازی
پشته در دو زبان را مقایسه می‌کند.

\begin{matintable}{مقایسه‌ی پیاده‌سازی پشته در پایتون و \lr{C++}}{tab:stack-comparison}
\begin{tblr}{
    width = \textwidth,
    colspec = {l X[c] X[c]},
    hlines, vlines,
    colsep = 8pt,
    rowsep = 4pt,
    row{1} = {font=\bfseries},
}
ویژگی & پایتون & \lr{C++} \\
نوع‌دهی & پویا & ایستا \\
ساختار پایه & \lr{list} & \lr{vector} \\
قالب & ندارد & دارد \\
مدیریت خطا & \lr{IndexError} & \lr{std::runtime\_error} \\
کارایی & متوسط & بالا \\
\end{tblr}
\end{matintable}

\mnote{در پایتون، نوع‌دهی پویا و در \lr{C++}، نوع‌دهی ایستا است. هر دو رویکرد مزایا و معایب خود را دارند.}

\section{تمرین‌ها}
\label{sec:programming-book-exercises}

\begin{exercise}
    یک تابع بنویسید که با استفاده از پشته،
    یک رشته را معکوس کند.
    \label{exc:reverse-string}
\end{exercise}

\begin{solution}
    \begin{latin}
\begin{minted}{python}
def reverse_string(s: str) -> str:
    """Reverse a string using a stack."""
    stack = Stack()
    for char in s:
        stack.push(char)

    result = []
    while not stack.is_empty():
        result.append(stack.pop())

    return "".join(result)


# |\pc{تست}|
print(reverse_string("hello"))  # "olleh"
\end{minted}
    \end{latin}
\end{solution}

\begin{exercise}
    با استفاده از پشته، یک ماشین حساب ساده
    برای عبارات پسوندی (\lr{postfix}) بنویسید.
    \label{exc:postfix-calculator}
\end{exercise}

\begin{exercise}
    تفاوت پشته و صف را توضیح دهید و یک
    کاربرد واقعی برای هر یک بنویسید.
    \label{exc:stack-vs-queue}
\end{exercise}

\section{خلاصه}
\label{sec:programming-book-summary}

در این فصل، یک فصل کامل از یک کتاب
برنامه‌نویسی را دیدیم که شامل:

\begin{itemize}
    \item \textbf{تعاریف}: پشته، \lr{LIFO}.
    \item \textbf{کد}: پیاده‌سازی در پایتون
    و \lr{C++}.
    \item \textbf{الگوریتم}: بررسی توازن
    پرانتزها.
    \item \textbf{قضیه}: پیچیدگی الگوریتم.
    \item \textbf{اثبات}: تحلیل پیچیدگی.
    \item \textbf{جدول}: مقایسه‌ی دو زبان.
    \item \textbf{تمرین‌ها}: با راه‌حل.
\end{itemize}

\mnote{این فصل، الگویی برای نوشتن فصل‌های کتاب برنامه‌نویسی شماست. می‌توانید ساختار آن را برای موضوعات دیگر (مثل درخت، گراف، مرتب‌سازی) تطبیق دهید.}

در فصل بعدی (\cref{chap:physics-book})،
یک فصل نمونه از یک کتاب فیزیک را می‌بینیم
که بر معادلات و نمودارهای فیزیکی تمرکز
دارد.